\documentclass[10pt]{article}
\usepackage[
  a5paper,
  top=2cm,
  bottom=2cm,
  left=2cm,
  right=2cm
]{geometry}
\usepackage[bookmarks]{hyperref}
\usepackage[T1]{fontenc}
\usepackage[utf8]{inputenc}
\usepackage{xpatch}
\usepackage[chorded,noshading,nomeasures]{songs}
% \usepackage[chorded,transposecapos,noshading,nomeasures]{songs}

% =================================================================
% TOGGLE SWITCH: Set to 1 to hide subsequent choruses, 0 to show all
\newcommand{\HideSubsequentChoruses}{0}
% =================================================================

\setlength{\parindent}{0pt}
\setlength{\sbarheight}{0pt}
\nosongnumbers{}

\makeatletter

\newcounter{choruscount}

% --- Change spacing ---
\versesep=14pt plus 6pt minus 8pt
\baselineadj=0pt plus 0pt minus 0pt

% === Patch chorus formating ===
\xpatchcmd{\justifyleft}
  {\SB@cbarshift}
  {\ifSB@inchorus\advance\leftskip\versenumwidth\fi}
  {}{}

\xpatchcmd{\justifycenter}
  {\SB@cbarshift}{}
  {}{}

\xpatchcmd{\SB@par}
  {\ifdim\cbarwidth>\z@\nobreak\else\SB@ilpenalty\fi}{}
  {}{}

\xpatchcmd{\brk}
  {
    \ifdim\cbarwidth=\z@%
    \ifrepchorus\marks\SB@cmarkclass{}\fi%
    \SB@breakpoint\brkpenalty%
  }
  {
    \endgroup\egroup%
    \ifrepchorus\ifSB@gotchorus\else
      \global\setbox\SB@chorusbox\vbox{
        \unvbox\SB@chorusbox%
        \SB@chorusbar\SB@box%
        \unvcopy\SB@box%
        \SB@breakpoint\brkpenalty%
      }%
    \fi
  }
  {}{}

\xpatchcmd{\beginchorus}
  {\vnumberedfalse}{\vnumberedtrue}
  {}{}

\xpretocmd{\beginchorus}{\stepcounter{choruscount}}{}{}

\xpretocmd{\endchorus}{%
  \ifnum\HideSubsequentChoruses=1\relax%
    \ifnum\value{choruscount}>1\relax%
      \endgroup\egroup%
    \fi%
  \fi%
}{}{}

\xpatchcmd{\SB@@beginchorus}
  {
    \setbox\SB@box\vbox\bgroup\begingroup%
      \ifchorded%
        \def\SB@everypar{%
          \vrule\@height\baselineskip\@width\z@\@depth\z@%
          \gdef\SB@everypar{}%
        }%
        \everypar{\SB@everypar\everypar{}}%
      \fi%
  }
  {
    \setbox\SB@box\vbox\bgroup\begingroup
      \ifvnumbered%
        \def\SB@everypar{%
          \ifnum\HideSubsequentChoruses=1\relax%
            \ifnum\value{choruscount}>1\relax%
              \chordedfalse%
              \leavevmode\hbox%
              to \versenumwidth{R:\hfil}%
              \gdef\SB@everypar{}%
              \setbox\z@\vbox\bgroup\begingroup%
              \chordedtrue%
            \else%
              \setbox\SB@box\hbox{R:}%
              \ifdim\wd\SB@box<\versenumwidth%
                \setbox\SB@box\hbox%
                to\versenumwidth{\unhbox\SB@box\hfil}%
              \fi%
              \ifchorded\vrule\@height\baselineskip\@width\z@\@depth\z@\fi%
              \placeversenum\SB@box%
              \gdef\SB@everypar{}%
            \fi%
          \else%
            \setbox\SB@box\hbox{R:}%
            \ifdim\wd\SB@box<\versenumwidth%
              \setbox\SB@box\hbox%
              to\versenumwidth{\unhbox\SB@box\hfil}%
            \fi%
            \ifchorded\vrule\@height\baselineskip\@width\z@\@depth\z@\fi%
            \placeversenum\SB@box%
            \gdef\SB@everypar{}%
          \fi%
        }
      \else
        \def\SB@everypar{
          \ifchorded\vrule\@height\baselineskip\@width\z@\@depth\z@\fi
          \gdef\SB@everypar{}
        }
      \fi
    \everypar{\SB@everypar\everypar{}}
  }
  {}{}

\xpatchcmd{\SB@endchorus}
  {\SB@chorusbar\SB@box}{}
  {}{}

\newcommand{\InputSongFolder}[1]{%
  \typeout{Scanning for songs inside #1\ldots}%
  \openin\@inputcheck=|'find #1 -name '*.tex' -type f | sort'%
  \loop%
    \read\@inputcheck%
    to \@songpath%
    \unless\ifeof\@inputcheck%
      \filename@parse{\@songpath}%
      \edef\cleanpath{\filename@area\filename@base.\filename@ext}%
      \expandafter\input\expandafter{\cleanpath}%
    \repeat%
  \closein\@inputcheck%
}

\newcommand{\allsongs}{}
\def\addsong#1{\g@addto@macro\allsongs{\input{src/#1}\sclearpage}}

\makeatother
% === Patch chorus formating ===

% === Patch Table of Contents ===
\makeatletter
\let\oldbeginsong\beginsong%

\renewcommand\beginsong[1]{%
  \setcounter{choruscount}{0}%
  \hypersetup{bookmarksdepth=0}%
  \phantomsection%
  \addcontentsline{toc}{subsection}{#1}%
  \hypersetup{bookmarksdepth=2}%
  \oldbeginsong{#1}%
}
\makeatother
% === Patch Table of Contents ===

\renewcommand{\contentsname}{Obash}

\begin{document}

\tableofcontents
\clearpage

\renewcommand{\lyricfont}{\small}
\renewcommand{\printchord}{\it\small\/}

\renewcommand{\capo}[1]{ % fix capo spacing
  \ifchorded{}
    \vspace{-1.5ex}\musicnote{Capo #1}\vspace{-5.5ex}
  \fi
}

\songcolumns{1}

% convert from engish notation to german notation
% \notenamesin{A}{B}{C}{D}{E}{F}{G}
% \notenamesout{A}{H}{C}{D}{E}{F}{G}

% \makeatletter
% \def\SB@bchord#1{%
%   \ifnum\SB@modifier=-\@ne\if\SB@basename{} B%
%     \def\SB@outstring{B}%
%     \SB@modifier=\z@%
%   \fi\fi%
%   \SB@chordfont\SB@outstring\futurelet\next\SB@echord%
% }
% \makeatother

% \renewcommand{\sharpsymbol}{is}
% \renewcommand{\flatsymbol}{es}
% ====================================================================

\begin{songs}{}
  \allsongs{
    % Generated automatically by convert.py (Musicians Layout Profile)
\input{songs/oddilovy_zpevnik_II/src/Amazonka-Hop_Trop+musician.tex}\sclearpage{}
\input{songs/oddilovy_zpevnik_II/src/Bláznova_ukolébavka+musician.tex}\sclearpage{}
% chktex-file 8
% chktex-file 12
% chktex-file 18
\beginsong{Cesta}[by={Kryštof}]
% \capo{0}
\transpose{0}

\beginverse*{}
\ch{C}{}{}{} \ch{G}{}{}{} \ch{Ami}{}{}{} \ch{F}{}{}{} \ch{G}{}{}{}
\endverse{}

\beginverse{}
Tou \ch{C}{}{}{}cestou, tím směrem prý bych se \ch{G}{}{}{}dávno měl dát,
když \ch{Dmi}{}{}{}sněží, jde to stěží, ale sněhy pak tají,
kus \ch{F}{}{}{}něhy ti za nehty \ch{G}{}{}{}slíbí a dají.
\endverse{}

\beginverse{}
Víc \ch{C}{}{}{}síly, se prát, na dně víc \ch{G}{}{}{}dávat než brát,
a i\ch{Dmi}{}{}{} když se vleče a je schůdná jen v kleče
donutí \ch{F}{}{}{}přestat se zbytečně \ch{G}{}{}{}ptát.
\endverse{}

\beginchorus{}
Jestli se \ch{C}{}{}{}blížím k cíli, kolik \ch{G}{}{}{}zbývá víry,
kam \ch{Ami}{}{}{}zvou, svodidla, co po tmě mi \ch{F}{}{}{}lžou? \ch{G}{}{}{}
Zda \ch{C}{}{}{}couvám zpátky, a \ch{G}{}{}{}plýtvám řádky, co \ch{Ami}{}{}{}řvou,
že už mi doma neotev\ch{F}{}{}{}řou. \ch{G}{}{}{}
\endchorus{}

\beginverse{}
Nebo jít s \ch{C}{}{}{}proudem, na lusknutí prstu se \ch{G}{}{}{}začít hned smát,
mít \ch{Dmi}{}{}{}svůj chodník slávy a před sebou davy
a přes \ch{F}{}{}{}zkroucená záda být \ch{G}{}{}{}součástí stáda.
\endverse{}

\beginverse{}
Ale \ch{C}{}{}{}zpívat, a hrát, kotníky \ch{G}{}{}{}líbat, a stát,
na \ch{Dmi}{}{}{}křídlech všech slavíků, a vlastně už ze zvyku
\ch{F}{}{}{}přestat se zbytečně \ch{G}{}{}{}ptát.
\endverse{}

\repeatingchorus{}
\beginchorus{}
Jestli se \ch{C}{}{}{}blížím k cíli, kolik \ch{G}{}{}{}zbývá víry,
kam \ch{Ami}{}{}{}zvou, svodidla, co po tmě mi \ch{F}{}{}{}lžou? \ch{G}{}{}{}
Zda \ch{C}{}{}{}couvám zpátky, a \ch{G}{}{}{}plýtvám řádky, co \ch{Ami}{}{}{}řvou,
že už mi doma neotev\ch{F}{}{}{}řou. \ch{G}{}{}{}
\endchorus{}

\repeatingchorus{}
\beginchorus{}
Jestli se \ch{C}{}{}{}blížím k cíli, kolik \ch{G}{}{}{}zbývá víry,
kam \ch{Ami}{}{}{}zvou, svodidla, co po tmě mi \ch{F}{}{}{}lžou? \ch{G}{}{}{}
Zda \ch{C}{}{}{}couvám zpátky, a \ch{G}{}{}{}plýtvám řádky, co \ch{Ami}{}{}{}řvou,
že už mi doma neotev\ch{F}{}{}{}řou. \ch{G}{}{}{}
\endchorus{}

\endsong{}
\sclearpage{}
\input{songs/oddilovy_zpevnik_II/src/Divokej_horskej_tymián-Žalman+musician.tex}\sclearpage{}
\input{songs/oddilovy_zpevnik_II/src/Gorale-Čechomor+musician.tex}\sclearpage{}
\input{songs/oddilovy_zpevnik_II/src/Hudsonský_šífy-Wabi_Daněk+musician.tex}\sclearpage{}
\input{songs/oddilovy_zpevnik_II/src/Husličky-Hradišťan+musician.tex}\sclearpage{}
% chktex-file 8
% chktex-file 12
% chktex-file 18
\beginsong{Jantarová země}[by={Žalman}]
% \capo{0}
\transpose{0}

\beginverse{}
\ch{Ami}{}{}{}Do splašenejch \ch{C}{}{}{}koní už zase lovci \ch{Ami}{}{}{}střílí
ulicí je \ch{C}{}{}{}ženou dolů k ohra\ch{Emi}{}{}{}dám
zítra jako v\ch{Dmi7}{}{}{}loni hra je pro ně \ch{Ami}{}{}{}dohraná
bujný hlavy sk\ch{Dmi7}{}{}{}loní někam do trá\ch{Ami}{}{}{}vy
\endverse{}

\beginchorus{}
\ch{Dmi7}{}{}{}Krá, až si kabát \ch{Ami}{}{}{}změní
\ch{Dmi7}{}{}{}krá, bílé vrány odle\ch{Ami}{}{}{}tí
\ch{Dmi7}{}{}{}krá, v Jantarový \ch{Ami}{}{}{}zemi
přista\ch{F}{}{}{}nou v jiném stole\ch{Emi}{}{}{}tí
\endchorus{}

\beginverse*{}
\ch{Em}{}{}{} \ch{Dmi}{}{}{} \ch{Ami}{}{}{} \ch{Dmi}{}{}{} \ch{Ami}{}{}{}
\endverse{}

\beginverse{}
\ch{Ami}{}{}{}Zamčený jsou \ch{C}{}{}{}stáje kat už hřívy \ch{Ami}{}{}{}stříhá
lovčí král je v \ch{C}{}{}{}kapli po stý koruno\ch{Emi}{}{}{}ván
sníh pomalu \ch{Dmi7}{}{}{}taje těla stromů se \ch{Ami}{}{}{}probouzí
v Jantarový \ch{Dmi7}{}{}{}zemi bude zas co \ch{Ami}{}{}{}pít
\endverse{}

\beginchorus{}
\ch{Dmi7}{}{}{}Krá, až si kabát \ch{Ami}{}{}{}změní
\ch{Dmi7}{}{}{}krá, bílé vrány odle\ch{Ami}{}{}{}tí
\ch{Dmi7}{}{}{}krá, v Jantarový \ch{Ami}{}{}{}zemi
přista\ch{F}{}{}{}nou v jiném stole\ch{Emi}{}{}{}tí
\endchorus{}

\beginverse*{}
\ch{Em}{}{}{} \ch{Dmi}{}{}{} \ch{Ami}{}{}{} \ch{Dmi}{}{}{} \ch{Ami}{}{}{}
\endverse{}

\endsong{}
\sclearpage{}
\input{songs/oddilovy_zpevnik_II/src/Jasná_zpráva-Olympic+musician.tex}\sclearpage{}
% chktex-file 8
% chktex-file 12
% chktex-file 18
\beginsong{Jdou po mně, jdou}[by={Jaromír Nohavica}]
% \capo{0}
\transpose{0}

\beginverse{}
Býval jsem \ch{D}{}{}{}chudý jak \ch{G}{}{}{}kostelní \ch{D}{}{}{}myš,
na půdě \ch{F{\shrp}mi}{}{}{}půdy jsem m\ch{Bmi}{}{}{}íval svou s\ch{A}{}{}{}krýš,
\lrep{} \ch{G}{}{}{}pak jednou v \ch{D}{}{}{}létě \ch{A}{}{}{}řek' jsem si: \ch{Bmi}{}{}{}bať,
\ch{G}{}{}{}svět facku\ch{D}{}{}{}je tě, a \ch{G}{}{}{}tak mu to v\ch{D}{}{}{}rať. \rrep{}
\endverse{}

\beginverse{}
Když mi dát \ch{D}{}{}{}nechceš, já \ch{G}{}{}{}vezmu si \ch{D}{}{}{}sám,
zámek jde \ch{F{\shrp}mi}{}{}{}lehce a \ch{Bmi}{}{}{}adresu \ch{A}{}{}{}znám,
\lrep{} \ch{G}{}{}{}zlato jak \ch{D}{}{}{}zlato, \ch{A}{}{}{}dolar či \ch{Bmi}{}{}{}frank,
\ch{G}{}{}{}tak jsem šel \ch{D}{}{}{}na to do \ch{G}{}{}{}National \ch{D}{}{}{}Bank. \rrep{}
\endverse{}

\beginchorus{}
Jdou po mně, \ch{D}{}{}{}jdou, \ch{G}{}{}{}jdou, \ch{D}{}{}{}jdou,
na každém \ch{F{\shrp}mi}{}{}{}rohu ma\ch{Bmi}{}{}{}jí fotku \ch{A}{}{}{}mou,
\ch{G}{}{}{}kdyby mě \ch{D}{}{}{}chytli, \ch{A}{}{}{}jó, byl by \ch{Bmi}{}{}{}ring,
\ch{G}{}{}{}tma jako v \ch{D}{}{}{}pytli je v \ch{C}{}{}{}celách Sing-\ch{G}{}{}{}sing. \ch{A}{}{}{}Jé, \ch{D}{}{}{}Jé.\ch{G}{}{}{}.. \ch{D}{}{}{}
\endchorus{}

\beginverse{}
Ve státě \ch{D}{}{}{}Iowa byl \ch{G}{}{}{}od poldů \ch{D}{}{}{}klid,
chudinká \ch{F{\shrp}mi}{}{}{}vdova mi \ch{Bmi}{}{}{}nabídla \ch{A}{}{}{}byt,
\lrep{} \ch{G}{}{}{}byla to \ch{D}{}{}{}kráska, \ch{A}{}{}{}já měl pení\ch{Bmi}{}{}{}ze,
\ch{G}{}{}{}tak začla \ch{D}{}{}{}láska jak \ch{G}{}{}{}z televi\ch{D}{}{}{}ze. \rrep{}
\endverse{}

\beginverse{}
Však půl roku \ch{D}{}{}{}nato \ch{G}{}{}{}řekla mi:"\ch{D}{}{}{}Dost,
tobě došlo \ch{F{\shrp}mi}{}{}{}zlato, mně \ch{Bmi}{}{}{}trpěli\ch{A}{}{}{}vost,
\lrep{} \ch{G}{}{}{}sbal svých pár \ch{D}{}{}{}švestek a \ch{A}{}{}{}běž si, kam \ch{Bmi}{}{}{}chceš,"
\ch{G}{}{}{}tak jsem na \ch{D}{}{}{}cestě a \ch{G}{}{}{}chudý jak \ch{D}{}{}{}veš. \rrep{}
\endverse{}

\beginchorus{}
Jdou po mně, \ch{D}{}{}{}jdou, \ch{G}{}{}{}jdou, \ch{D}{}{}{}jdou,
na každém \ch{F{\shrp}mi}{}{}{}rohu ma\ch{Bmi}{}{}{}jí fotku \ch{A}{}{}{}mou,
\ch{G}{}{}{}kdyby mě \ch{D}{}{}{}chytli, \ch{A}{}{}{}jó, byl by \ch{Bmi}{}{}{}ring,
\ch{G}{}{}{}tma jako v \ch{D}{}{}{}pytli je v \ch{C}{}{}{}celách Sing-\ch{G}{}{}{}sing. \ch{A}{}{}{}Jé, \ch{D}{}{}{}Jé.\ch{G}{}{}{}.. \ch{D}{}{}{}
\endchorus{}

\beginverse{}
Teď ve státě \ch{D}{}{}{}Utah ži\ch{G}{}{}{}ju spoko\ch{D}{}{}{}jen,
pípu jsem \ch{F{\shrp}mi}{}{}{}utáh' a \ch{Bmi}{}{}{}straním se \ch{A}{}{}{}žen,
\lrep{} \ch{G}{}{}{}kladou mi \ch{D}{}{}{}pasti a \ch{A}{}{}{}do pastí \ch{Bmi}{}{}{}špek,
\ch{G}{}{}{}já na ně \ch{D}{}{}{}mastím, jen \ch{G}{}{}{}ať mají \ch{D}{}{}{}vztek. \rrep{}
\endverse{}

\beginchorus{}
Jdou po mně \ch{D}{}{}{}jdou, \ch{G}{}{}{}jdou, \ch{D}{}{}{}jdou,
na nočních \ch{F{\shrp}mi}{}{}{}stolcích ma\ch{Bmi}{}{}{}jí fotku \ch{A}{}{}{}mou,
\ch{G}{}{}{}kdyby mě \ch{D}{}{}{}klofly, \ch{A}{}{}{}jó, byl by \ch{Bmi}{}{}{}ring,
\ch{G}{}{}{}žít pod pan\ch{D}{}{}{}toflí je \ch{C}{}{}{}hůř než v Sing-\ch{G}{}{}{}sing. \ch{A}{}{}{}Jé, \ch{D}{}{}{}Jé.\ch{G}{}{}{}.. \ch{D}{}{}{}
\endchorus{}

\endsong{}
\sclearpage{}
% chktex-file 8
% chktex-file 12
% chktex-file 18
\beginsong{Jelen}[by={Jelen}]
% \capo{0}
\transpose{0}

\beginverse{}
\ch{Dmi}{}{}{}Na jaře se vrací, \ch{C}{}{}{}od podzima lis\ch{Dmi}{}{}{}tí.
Mraky místo ptáků, \ch{C}{}{}{}krouží nad závis\ch{Dmi}{}{}{}tí.
Kdyby jsi se někdy \ch{C}{}{}{}ke mně chtěla vrá\ch{Dmi}{}{}{}tit,
nesměla bys, lásko, \ch{C}{}{}{}moje srdce ztra\ch{G}{}{}{}tit.
\endverse{}

\beginchorus{}
\ch{Dmi}{}{}{}Zabil jsem v lese \ch{C}{}{}{}jele\ch{F}{}{}{}na,
bez nenávisti, \ch{C}{}{}{}bez jmé\ch{Dmi}{}{}{}na.
Když přišel dolů k \ch{C}{}{}{}řece \ch{F}{}{}{}pít,
krev teče do vody, \ch{C}{}{}{}v srdci \ch{Dmi}{}{}{}klid.
\endchorus{}

\beginverse{}
\ch{Dmi}{}{}{}Voda teče k moři, \ch{C}{}{}{}po kamenech ská\ch{Dmi}{}{}{}če,
jednou hráze boří, \ch{C}{}{}{}jindy tiše plá\ch{Dmi}{}{}{}če.
Někdy mám ten pocit, \ch{C}{}{}{}ikdyž roky ply\ch{Dmi}{}{}{}nou,
že vidím tvůj odraz, \ch{C}{}{}{}dole pod hladi\ch{G}{}{}{}nou.
\endverse{}

\beginchorus{}
\ch{Dmi}{}{}{}Zabil jsem v lese \ch{C}{}{}{}jele\ch{F}{}{}{}na,
bez nenávisti, \ch{C}{}{}{}bez jmé\ch{Dmi}{}{}{}na.
Když přišel dolů k \ch{C}{}{}{}řece \ch{F}{}{}{}pít,
krev teče do vody, \ch{C}{}{}{}v srdci \ch{Dmi}{}{}{}klid.
\endchorus{}

\beginverse{}
\ch{Dmi}{}{}{}Na jaře se vrací, \ch{C}{}{}{}listí od podzi\ch{Dmi}{}{}{}ma.
Čas se někam ztrácí, \ch{C}{}{}{}brzo bude zi\ch{Dmi}{}{}{}ma.
Svět přikryje ticho, \ch{C}{}{}{}tečka za příbě\ch{Dmi}{}{}{}hem.
Kdo pozná, čí kosti, \ch{C}{}{}{}zapadaly sně\ch{G}{}{}{}hem.
\endverse{}

\beginchorus{}
\ch{Dmi}{}{}{}Zabil jsem v lese \ch{C}{}{}{}jele\ch{F}{}{}{}na,
bez nenávisti, \ch{C}{}{}{}bez jmé\ch{Dmi}{}{}{}na.
Když přišel dolů k \ch{C}{}{}{}řece \ch{F}{}{}{}pít,
krev teče do vody, \ch{C}{}{}{}v srdci \ch{Dmi}{}{}{}klid.
\endchorus{}

\endsong{}
\sclearpage{}
\input{songs/oddilovy_zpevnik_II/src/John_Brown+musician.tex}\sclearpage{}
\input{songs/oddilovy_zpevnik_II/src/Joy-Kabát+musician.tex}\sclearpage{}
% chktex-file 8
% chktex-file 12
% chktex-file 18
\beginsong{Kdyby se v komnatách}[by={Zdeněk Svěrák, Jaroslav Uhlíř}]
% \capo{0}
\transpose{0}

\beginverse{}
\ch{G}{}{}{}K životu na zámku, \ch{Dm}{}{}{}mám jednu poznámku,
\ch{C}{}{}{}je tu neveselo, je tu \ch{G}{}{}{}truchlivo.
V ostatních královstvích, \ch{Dm}{}{}{}nezní tak málo smích,
\ch{C}{}{}{}není neveselo, není \ch{G}{}{}{}truchlivo.
\endverse{}

\beginchorus{}
\ch{G7}{}{}{}Kdyby se v \ch{C}{}{}{}komnatách, běhoun jak \ch{Bm}{}{}{}hrom natáh,
a na něm \ch{Ami}{}{}{}akrobati,\ch{D}{}{}{} začali \ch{G}{}{}{}kejklovati.
\ch{G7}{}{}{}Kdyby nám v \ch{C}{}{}{}paláci, pištěli \ch{Bm}{}{}{}dudáci,
to by se \ch{Ami}{}{}{}krásně žilo,\ch{D}{}{}{} to by byl \ch{G}{}{}{}ráj.
\endchorus{}

\beginverse{}
\ch{G}{}{}{}Kde není muzika, \ch{Dm}{}{}{}tam duše naříká,
\ch{C}{}{}{}tam je neveselo, tam je \ch{G}{}{}{}truchlivo.
\ch{G}{}{}{}Chtěla bych dvůr pestrý, \ch{Dm}{}{}{}kde znějí orchestry,
\ch{C}{}{}{}pěkně na veselo, žádné \ch{G}{}{}{}truchlivo.
\endverse{}

\beginchorus{}
\ch{G7}{}{}{}Kdyby se v \ch{C}{}{}{}komnatách, běhoun jak \ch{Bm}{}{}{}hrom natáh,
a na něm \ch{Ami}{}{}{}akrobati,\ch{D}{}{}{} začali \ch{G}{}{}{}kejklovati.
\ch{G7}{}{}{}Kdyby nám v \ch{C}{}{}{}paláci, pištěli \ch{Bm}{}{}{}dudáci,
to by se \ch{Ami}{}{}{}krásně žilo,\ch{D}{}{}{} to by byl \ch{G}{}{}{}ráj.
\endchorus{}

\endsong{}
\sclearpage{}
\input{songs/oddilovy_zpevnik_II/src/Když_mě_brali_za_vojáka-Jaromír_Nohavica+musician.tex}\sclearpage{}
% chktex-file 8
% chktex-file 12
% chktex-file 18
\beginsong{Kláda}[by={Hop Trop}]
% \capo{0}
\transpose{0}

\beginverse{}
\ch{Ami}{}{}{}Celý roky prachy jsem si skládal,
\ch{C}{}{}{}nikdy svýho f\ch{G}{}{}{}loka nepro\ch{Ami}{}{}{}pil,
\ch{Ami}{}{}{}vod lopaty měl vohnutý záda,
\ch{C}{}{}{}paty od baráku \ch{G}{}{}{}jsem neodle\ch{Ami}{}{}{}pil,
\ch{G}{}{}{}nikdo neví, do čeho se \ch{Ami}{}{}{}někdy zamotá,
tohle \ch{C}{}{}{}já už \ch{G}{}{}{}dávno pocho\ch{Ami}{}{}{}pil.
\endverse{}

\beginverse{}
\ch{Ami}{}{}{}Taky kdysi vydělat jsem toužil,
\ch{C}{}{}{}brácha řek' mi, \ch{G}{}{}{}že by se mnou \ch{Ami}{}{}{}šel,
\ch{Ami}{}{}{}tak jsem háky, lana, klíny koupil,
\ch{C}{}{}{}a sekyru jsme svoji \ch{G}{}{}{}každej doma \ch{Ami}{}{}{}měl,
\ch{G}{}{}{}a plány veliký, jak \ch{Ami}{}{}{}fajn budem se mít,
nikdo z \ch{C}{}{}{}nás pro \ch{G}{}{}{}holku nebre\ch{Ami}{}{}{}čel.
\endverse{}

\beginchorus{}
\ch{F}{}{}{}Duní \ch{C}{}{}{}kláda kory\ch{Dmi}{}{}{}tem, bacha \ch{E}{}{}{}dej, hej, bacha \ch{Ami}{}{}{}dej!
S \ch{F}{}{}{}tou si, \ch{C}{}{}{}bráško, nety\ch{Dmi}{}{}{}kej, nety\ch{E}{}{}{}kej!
\endchorus{}

\beginverse{}
\ch{Ami}{}{}{}Dřevo dostat k pile, kde ho koupí,
\ch{C}{}{}{}není těžký, \ch{G}{}{}{}vždyť jsme \ch{Ami}{}{}{}fikaný,
\ch{Ami}{}{}{}ten rok bylo jaro ale skoupý,
\ch{C}{}{}{}a teď jsme na dně my i \ch{G}{}{}{}vory sváza\ch{Ami}{}{}{}ný,
\ch{G}{}{}{}a k tomu můžem říct jen, \ch{Ami}{}{}{}že nemáme nic,
jen kus \ch{C}{}{}{}práce \ch{G}{}{}{}nedoděla\ch{Ami}{}{}{}ný.
\endverse{}

\beginchorus{}
\lrep{} \ch{F}{}{}{}Duní \ch{C}{}{}{}kláda kory\ch{Dmi}{}{}{}tem, bacha \ch{E}{}{}{}dej, hej, bacha \ch{Ami}{}{}{}dej!
S \ch{F}{}{}{}tou si, \ch{C}{}{}{}bráško, nety\ch{Dmi}{}{}{}kej, nety\ch{E}{}{}{}kej! \rrep{} \ch{Ami}{}{}{}
\endchorus{}

\endsong{}
\sclearpage{}
% chktex-file 8
% chktex-file 12
% chktex-file 18
\beginsong{Kluziště}[by={Karel Plíhal}]
% \capo{0}
\transpose{0}

\beginverse{}
\ch{C}{}{}{}Strejček \ch{Emi}{}{}{}kovář, \ch{Ami}{}{}{}chytil \ch{C}{}{}{}kleště, \ch{Fmaj7}{}{}{}uštíp' z \ch{C}{}{}{}noční \ch{Fmaj7}{}{}{}oblohy\ch{G}{}{}{}
jednu malou, kapku deště, ta mu spadla pod nohy.
Nejdřív ale chytil slinu, tak sáh' kamsi pro pivo,
pak přitáhnul, kovadlinu a obrovský kladivo.
\endverse{}

\beginchorus{}
Zatím \ch{C}{}{}{}tři bílé \ch{Emi}{}{}{}vrány \ch{Ami}{}{}{}pěkně za se\ch{C}{}{}{}bou,
kolem \ch{Fmaj7}{}{}{}jdou, někam \ch{C}{}{}{}jdou, do ry\ch{D7}{}{}{}tmu se kýva\ch{G}{}{}{}jí.
Tyhle \ch{C}{}{}{}tři bílé \ch{Emi}{}{}{}vrány \ch{Ami}{}{}{}pěkně za se\ch{C}{}{}{}bou,
kolem \ch{Fmaj7}{}{}{}jdou, někam \ch{C}{}{}{}jdou, nedo\ch{Fmaj7}{}{}{}jdou, nedo\ch{C}{}{}{}jdou.
\endchorus{}

\beginverse{}
\ch{C}{}{}{}Vydal z \ch{Emi}{}{}{}hrdla, \ch{Ami}{}{}{}mocný \ch{C}{}{}{}pokřik, \ch{Fmaj7}{}{}{}ztichlým \ch{C}{}{}{}letním \ch{Fmaj7}{}{}{}večere\ch{G}{}{}{}m,
pak tu kapku, všude rozstřík', jedním mocným úderem.
Celej svět byl, náhle v kapce a vysoko nad námi,
na obrovské mucholapce, visí nebe s hvězdami.
\endverse{}

\repeatingchorus{}
\beginchorus{}
Zatím \ch{C}{}{}{}tři bílé \ch{Emi}{}{}{}vrány \ch{Ami}{}{}{}pěkně za se\ch{C}{}{}{}bou,
kolem \ch{Fmaj7}{}{}{}jdou, někam \ch{C}{}{}{}jdou, do ry\ch{D7}{}{}{}tmu se kýva\ch{G}{}{}{}jí.
Tyhle \ch{C}{}{}{}tři bílé \ch{Emi}{}{}{}vrány \ch{Ami}{}{}{}pěkně za se\ch{C}{}{}{}bou,
kolem \ch{Fmaj7}{}{}{}jdou, někam \ch{C}{}{}{}jdou, nedo\ch{Fmaj7}{}{}{}jdou, nedo\ch{C}{}{}{}jdou.
\endchorus{}

\beginverse{}
\ch{C}{}{}{}Zpod ví\ch{Emi}{}{}{}ček mi, \ch{Ami}{}{}{}vytrysk' \ch{C}{}{}{}pramen, \ch{Fmaj7}{}{}{}na zmač\ch{C}{}{}{}kané \ch{Fmaj7}{}{}{}polštá\ch{G}{}{}{}ře,
kdosi mě vzal, kolem ramen a políbil na tváře.
Kdesi v dálce, rozmazaně, strejda kovář odchází,
do kalhot si čistí dlaně, umazané od sazí.
\endverse{}

\repeatingchorus{}
\beginchorus{}
Zatím \ch{C}{}{}{}tři bílé \ch{Emi}{}{}{}vrány \ch{Ami}{}{}{}pěkně za se\ch{C}{}{}{}bou,
kolem \ch{Fmaj7}{}{}{}jdou, někam \ch{C}{}{}{}jdou, do ry\ch{D7}{}{}{}tmu se kýva\ch{G}{}{}{}jí.
Tyhle \ch{C}{}{}{}tři bílé \ch{Emi}{}{}{}vrány \ch{Ami}{}{}{}pěkně za se\ch{C}{}{}{}bou,
kolem \ch{Fmaj7}{}{}{}jdou, někam \ch{C}{}{}{}jdou, nedo\ch{Fmaj7}{}{}{}jdou, nedo\ch{C}{}{}{}jdou.
\endchorus{}

\endsong{}
\sclearpage{}
\input{songs/oddilovy_zpevnik_II/src/Královna_z_Dundrum_Bay-Nedvědi+musician.tex}\sclearpage{}
% chktex-file 8
% chktex-file 12
% chktex-file 18
\beginsong{Královské reggae}
% \capo{0}
\transpose{0}

\beginverse*{}
\ch{D}{}{}{}Umbabau, umbabaumbaumba, \ch{Bmi}{}{}{}umbabau, umbabaumbaumba
\ch{D}{}{}{}Umbabau, umbabaumbaumba, \ch{Bmi}{}{}{}umbabau, umbabaumbaumba
\endverse{}

\beginverse{}
\ch{D}{}{}{}Někdo ráčí kráčet v brnění, má na něm rádoby \ch{A7}{}{}{}erb.
\ch{D}{}{}{}Když není urozen, \ch{G}{}{}{}rychle je prozrazen,
\ch{D}{}{}{}když ne\ch{A7}{}{}{}říká k\ch{D}{}{}{}rásné er, tak ten si \ch{A7}{}{}{}říká o \ch{D}{}{}{}malér.
\endverse{}

\beginverse{}
\ch{D}{}{}{}Moudrý král princeznin bratranec, když vdával svých pět \ch{A7}{}{}{}dcer.
\ch{D}{}{}{}Řek: "To dál nesnesu" \ch{G}{}{}{}a prchnul do lesů,
\ch{D}{}{}{}žádný \ch{A7}{}{}{}ženich \ch{D}{}{}{}nemá er, kdo nemá \ch{A7}{}{}{}er, ztrácí \ch{D}{}{}{}charakter.
\endverse{}

\beginchorus{}
\ch{G}{}{}{}Královské reggae, \ch{D}{}{}{}královské reggae,
\ch{G}{}{}{}královské reggae, \ch{A7}{}{}{}královské reggae,
\endchorus{}

\beginverse*{}
\ch{D}{}{}{}Umbabau, umbabaumbaumba, \ch{Bmi}{}{}{}umbabau, umbabaumbaumba
\ch{D}{}{}{}Umbabau, umbabaumbaumba, \ch{Bmi}{}{}{}umbabau, umbabaumbaumba
\endverse{}

\beginverse{}
\ch{D}{}{}{}Ten, kdo chce mít pravou princeznu, musí být z vyšších s\ch{A7}{}{}{}fér.
\ch{D}{}{}{}Kdo prosí o ruku, \ch{G}{}{}{}musí mít záruku, \ch{D}{}{}{}tou je \ch{A7}{}{}{}jeho \ch{D}{}{}{}ryčné er.
\ch{D}{}{}{}Pouze ten \ch{A7}{}{}{}může být \ch{D}{}{}{}premiér a kdo nemá \ch{A7}{}{}{}er, tak je \ch{D}{}{}{}amatér.
\endverse{}

\beginverse{}
\ch{D}{}{}{}Mějte, prosím, trochu strpení, rozhodně všechno je \ch{A7}{}{}{}fér.
\ch{D}{}{}{}Tak proč ty trampoty, \ch{G}{}{}{}příšerné drahoty,
\ch{D}{}{}{}přece \ch{A7}{}{}{}slyšíš \ch{D}{}{}{}to mé er, které mě \ch{A7}{}{}{}pasovalo nad prů\ch{D}{}{}{}měr.
\endverse{}

\repeatingchorus{}
\beginchorus{}
\ch{G}{}{}{}Královské reggae, \ch{D}{}{}{}královské reggae,
\ch{G}{}{}{}královské reggae, \ch{A7}{}{}{}královské reggae,
\endchorus{}

\beginverse*{}
\ch{D}{}{}{}Umbabau, umbabaumbaumba, \ch{Bmi}{}{}{}umbabau, umbabaumbaumba
\ch{D}{}{}{}Umbabau, umbabaumbaumba, \ch{Bmi}{}{}{}umbabau, umbabaumbaumba
\endverse{}

\endsong{}
\sclearpage{}
% chktex-file 8
% chktex-file 12
% chktex-file 18
\beginsong{Lachtani}[by={Jaromír Nohavica}]
% \capo{0}
\transpose{0}

\beginchorus{}
\lrep{} \ch{C}{}{}{}Lach, lach, \ch{F}{}{}{}jé, \ch{C}{}{}{}jé \ch{Ami}{}{}{}lach, lach, \ch{G}{}{}{}jé,\ch{C}{}{}{}jé \rrep{}
\endchorus{}

\beginverse{}
\ch{C}{}{}{}Jedna lachtaní \ch{F}{}{}{}rodi\ch{C}{}{}{}na
\ch{Ami}{}{}{}rozhodla se, že si vyjde \ch{G}{}{}{}do ki\ch{C}{}{}{}na,
\ch{C}{}{}{}jeli vlakem, metrem, lodí a pak \ch{F}{}{}{}tramva\ch{C}{}{}{}jí
a teď \ch{Ami}{}{}{}u kina Vesmír \ch{G}{}{}{}lachta\ch{C}{}{}{}jí,
\ch{G}{}{}{}lachtaní úspory \ch{C}{}{}{}dali dohromady,
\ch{G}{}{}{}koupili si lístky \ch{C}{}{}{}do první \ch{G}{}{}{}řady,
\ch{C}{}{}{}táta lachtan řekl: "Nebudem \ch{F}{}{}{}třít bí\ch{C}{}{}{}du"
a pro \ch{Ami}{}{}{}každého koupil pytlík \ch{G}{}{}{}araši\ch{C}{}{}{}dů.
\endverse{}

\repeatingchorus{}
\beginchorus{}
\lrep{} \ch{C}{}{}{}Lach, lach, \ch{F}{}{}{}jé, \ch{C}{}{}{}jé \ch{Ami}{}{}{}lach, lach, \ch{G}{}{}{}jé,\ch{C}{}{}{}jé \rrep{}
\endchorus{}

\beginverse{}
\ch{C}{}{}{}Na jižním pólu je \ch{F}{}{}{}nehez\ch{C}{}{}{}ky
a tak \ch{Ami}{}{}{}lachtani si vyjeli na \ch{G}{}{}{}grotes\ch{C}{}{}{}ky,
\ch{C}{}{}{}těšili se, jak bude \ch{F}{}{}{}vese\ch{C}{}{}{}lo,
když \ch{Ami}{}{}{}zazněl gong a v sále se \ch{G}{}{}{}setmě\ch{C}{}{}{}lo,
\ch{G}{}{}{}co to ale vidí jejich \ch{C}{}{}{}lachtaní zraky:
\ch{G}{}{}{}sníh a mráz a \ch{C}{}{}{}sněhové \ch{G}{}{}{}mraky,
\ch{C}{}{}{}pro veliký úspěch změna \ch{F}{}{}{}progra\ch{C}{}{}{}mu,
dnes \ch{Ami}{}{}{}dáváme film ze života \ch{G}{}{}{}lachta\ch{C}{}{}{}nů.
\endverse{}

\repeatingchorus{}
\beginchorus{}
\lrep{} \ch{C}{}{}{}Lach, lach, \ch{F}{}{}{}jé, \ch{C}{}{}{}jé \ch{Ami}{}{}{}lach, lach, \ch{G}{}{}{}jé,\ch{C}{}{}{}jé \rrep{}
\endchorus{}

\beginverse{}
\ch{C}{}{}{}Táta lachtan vyskočil ze \ch{F}{}{}{}sedad\ch{C}{}{}{}la
\ch{Ami}{}{}{}nevídaná zlost ho \ch{G}{}{}{}popad\ch{C}{}{}{}la
"\ch{C}{}{}{}Proto jsem se netrmácel \ch{F}{}{}{}přes celý \ch{C}{}{}{}svět,
\ch{Ami}{}{}{}abych tady v kině mrznul jako \ch{G}{}{}{}turecký \ch{C}{}{}{}med,
\ch{G}{}{}{}tady zima, doma zima, \ch{C}{}{}{}všude jen chlad,
\ch{G}{}{}{}kde má chudák lachtan \ch{C}{}{}{}relaxo\ch{G}{}{}{}vat?"
\ch{C}{}{}{}Nedivte se té lachtaní \ch{F}{}{}{}rodi\ch{C}{}{}{}ně,
že pak \ch{Ami}{}{}{}rozšlapala arašídy \ch{G}{}{}{}po ki\ch{C}{}{}{}ně.
\endverse{}

\repeatingchorus{}
\beginchorus{}
\lrep{} \ch{C}{}{}{}Lach, lach, \ch{F}{}{}{}jé, \ch{C}{}{}{}jé \ch{Ami}{}{}{}lach, lach, \ch{G}{}{}{}jé,\ch{C}{}{}{}jé \rrep{}
\endchorus{}

\beginverse{}
T\ch{C}{}{}{}ahle lachtaní \ch{F}{}{}{}rodi\ch{C}{}{}{}na
\ch{Ami}{}{}{}od té doby nechodí už \ch{G}{}{}{}do kina, lach, \ch{C}{}{}{}lach.
\endverse{}

\endsong{}
\sclearpage{}
\input{songs/oddilovy_zpevnik_II/src/Lajlaj_důdap_tabidybam-Bob_Fliedr+musician.tex}\sclearpage{}
\input{songs/oddilovy_zpevnik_II/src/Louisiana-Hop_Trop+musician.tex}\sclearpage{}
\input{songs/oddilovy_zpevnik_II/src/Lou_z_Lille-Klíč+musician.tex}\sclearpage{}
% chktex-file 8
% chktex-file 12
% chktex-file 18
\beginsong{Malování}[by={Divokej Bill}]
% \capo{0}
\transpose{0}

\beginverse{}
\ch{Emi}{}{}{}Nesnaž se, znáš se, \ch{C}{}{}{}řekni mi \ch{D}{}{}{}co je ji\ch{Emi}{}{}{}ný,
jak v kleci máš se \ch{C}{}{}{} \ch{D}{}{}{}pro nevin\ch{Emi}{}{}{}ný
noci dlouhý \ch{C}{}{}{}jsou plný \ch{D}{}{}{}touhy a \ch{Emi}{}{}{}lásky nás dvou \ch{C}{}{}{}
\endverse{}

\beginverse{}
\ch{D}{}{}{}Všechno hezký \ch{Emi}{}{}{}za sebou mám, \ch{C}{}{}{}můžu \ch{D}{}{}{}si za to \ch{Emi}{}{}{}sám, \ch{C}{}{}{}v hlavě \ch{D}{}{}{}hlavolam,
\ch{Emi}{}{}{}jen táta a máma \ch{C}{}{}{}jsou s \ch{D}{}{}{}náma, \ch{Emi}{}{}{}napořád s náma
\endverse{}

\beginchorus{}
\ch{C}{}{}{}To \ch{D}{}{}{}je to tvoje \ch{Emi}{}{}{}malování \ch{C}{}{}{}vzdušnejch \ch{D}{}{}{}zámků,
\ch{Emi}{}{}{}malování po \ch{C}{}{}{}zdech holejma \ch{D}{}{}{}rukama
tě \ch{Emi}{}{}{}nezachrání, už \ch{C}{}{}{}máš na ka\ch{D}{}{}{}hánku,
\ch{Emi}{}{}{}nezachrání, už jsi \ch{C}{}{}{}na zá\ch{D}{}{}{}dech,
je to \ch{Emi}{}{}{}za náma, ty čteš \ch{C}{}{}{}poslední \ch{D}{}{}{}stránku,
\ch{Emi}{}{}{}za náma, \ch{C}{}{}{} na \ch{D}{}{}{}zádech, \ch{Emi}{}{}{}za náma,
už \ch{C}{}{}{}máš na ka\ch{D}{}{}{}hánku, \ch{Emi}{}{}{}mezi náma, \ch{C}{}{}{}mi taky \ch{Ami}{}{}{}došel \ch{Bmi}{}{}{}dech.
\endchorus{}

\beginverse{}
\ch{Emi}{}{}{}Znáš se, \ch{C}{}{}{}řekni mi \ch{D}{}{}{}co je ji\ch{Emi}{}{}{}ný,
jak v kleci máš se, \ch{C}{}{}{} \ch{D}{}{}{}pro nevin\ch{Emi}{}{}{}ný
noci dlouhý \ch{C}{}{}{}jsou plný \ch{D}{}{}{}touhy a \ch{Emi}{}{}{}lásky nás tří.
\endverse{}

\endsong{}
\sclearpage{}
% chktex-file 8
% chktex-file 12
% chktex-file 18
\beginsong{Marie}[by={Tomáš Klus}]
% \capo{0}
\transpose{0}

\beginverse{}
Je \ch{C}{}{}{}den. Tak \ch{E}{}{}{}pojď Marie ven.
Budeme \ch{F}{}{}{}žít. A házet \ch{G}{}{}{}šutry do oken.
Je \ch{C}{}{}{}dva. Necháme \ch{E}{}{}{}doma trucovat.
Když nechtě\ch{F}{}{}{}jí -- nemusí. Nebu\ch{G}{}{}{}dem se vnucovat.
Jémi\ch{C}{}{}{}ne. Všechno \ch{E}{}{}{}zlý jednou pomine.
Tak \ch{F}{}{}{}Marie.\ch{G}{}{}{} Co ti je?
\endverse{}

\beginverse{}
Všemoc\ch{C}{}{}{}né. Jsou \ch{E}{}{}{}loutkařovi prsty.
Ať jsou \ch{F}{}{}{}tenký nebo tlustý. Občas \ch{G}{}{}{}přetrhají nit.
A to pak \ch{C}{}{}{}jít. A nemít \ch{E}{}{}{}nad sebou svý jistý.
Pořád s \ch{F}{}{}{}tváří optimisty. Listy v \ch{G}{}{}{}žití obracet.
\endverse{}

\beginverse{}
Je to \ch{C}{}{}{}jed. Mazat si \ch{E}{}{}{}kolem huby med.
A nesly\ch{F}{}{}{}šet. Jak se ti \ch{G}{}{}{}bortí svět.
Ma\ch{C}{}{}{}rie. Kdo přeží\ch{E}{}{}{}vá nežije. Tak ádi\ch{F}{}{}{}jé. \ch{G}{}{}{}
\endverse{}

\beginverse{}
\ch{C}{}{}{}Marie. Už zase \ch{E}{}{}{}máš (k)tulení sklony.
Jako \ch{F}{}{}{}loni. Slyším \ch{G}{}{}{}kostelní zvony znít.
\ch{C}{}{}{}A to mě zabije. \ch{E}{}{}{}A to mě zabije.
\ch{F}{}{}{}A to mě zabije. \ch{G}{}{}{}Jistojistě.
\endverse{}

\beginchorus{}
Já \ch{C}{}{}{}mám Marie \ch{E}{}{}{}rád. Když má \ch{F}{}{}{}moje bytí \ch{G}{}{}{}spád.
Býti \ch{C}{}{}{}věčně na ces\ch{E}{}{}{}tách. A k ránu \ch{F}{}{}{}spícím plícím.
\ch{G}{}{}{}Život vdechovat. \ch{C}{}{}{}Nechtěj mě milovat.
\ch{E}{}{}{}Nechtěj mě milovat. \ch{F}{}{}{}Nechtěj mě milovat. \ch{G}{}{}{}
\endchorus{}

\beginchorus{}
Já \ch{C}{}{}{}mám Marie \ch{E}{}{}{}rád. Když má \ch{F}{}{}{}moje bytí \ch{G}{}{}{}spád.
Býti \ch{C}{}{}{}věčně na ces\ch{E}{}{}{}tách. A k ránu \ch{F}{}{}{}spícím plícím.
\ch{G}{}{}{}Život vdechovat. \ch{C}{}{}{}Nechtěj mě milovat.
\ch{E}{}{}{}Nechtěj mě milovat. \ch{F}{}{}{}Nechtěj mě milovat. \ch{G}{}{}{}
\endchorus{}

\beginverse{}
\ch{C}{}{}{}Copak nemůže být. \ch{E}{}{}{}Mezi ženou a mužem.
\ch{F}{}{}{}Přátelství -- kde není \ch{G}{}{}{}nikdo nic dlužen. Prostě.
\ch{C}{}{}{}Jen prosté. \ch{E}{}{}{}Spříznění duší.
\ch{F}{}{}{}Aniž by kdokoli. \ch{G}{}{}{}Cokoli tušil.
\endverse{}

\beginverse{}
\ch{C}{}{}{}na na na.\ch{E}{}{}{}..\ch{F}{}{}{} \ch{G}{}{}{} \ch{C}{}{}{} \ch{E}{}{}{} \ch{F}{}{}{} \ch{G}{}{}{}
\endverse{}

\beginchorus{}
Já \ch{C}{}{}{}mám Marie \ch{E}{}{}{}rád. Když má \ch{F}{}{}{}moje bytí \ch{G}{}{}{}spád.
Býti \ch{C}{}{}{}věčně na ces\ch{E}{}{}{}tách. A k ránu \ch{F}{}{}{}spícím plícím.
\ch{G}{}{}{}Život vdechovat. \ch{C}{}{}{}Nechtěj mě milovat.
\ch{E}{}{}{}Nechtěj mě milovat. \ch{F}{}{}{}Nechtěj mě milovat. \ch{G}{}{}{}
\endchorus{}

\beginchorus{}
Já \ch{C}{}{}{}mám Marie \ch{E}{}{}{}rád. Když má \ch{F}{}{}{}moje bytí \ch{G}{}{}{}spád.
Býti \ch{C}{}{}{}věčně na ces\ch{E}{}{}{}tách. A k ránu \ch{F}{}{}{}spícím plícím.
\ch{G}{}{}{}Život vdechovat. \ch{C}{}{}{}Nechtěj mě milovat.
\ch{E}{}{}{}Nechtěj mě milovat. \ch{F}{}{}{}Nechtěj mě milovat. \ch{G}{}{}{}
\endchorus{}

\endsong{}
\sclearpage{}
\input{songs/oddilovy_zpevnik_II/src/Maruška-Malomocnost_Prázdnoty+musician.tex}\sclearpage{}
\input{songs/oddilovy_zpevnik_II/src/Měsíc-Mňága_&_Žďorp+musician.tex}\sclearpage{}
\input{songs/oddilovy_zpevnik_II/src/Mlýny-Spirituál_kvintet+musician.tex}\sclearpage{}
\input{songs/oddilovy_zpevnik_II/src/Mravenčí_ukolébavka-Zdeněk_Svěrák,_Jaroslav_Uhlíř+musician.tex}\sclearpage{}
\input{songs/oddilovy_zpevnik_II/src/Muži_z_Ria-Hana_Pazeltová+musician.tex}\sclearpage{}
% chktex-file 8
% chktex-file 12
% chktex-file 18
\beginsong{Nad stádem koní}[by={Buty}]
% \capo{0}
\transpose{0}

\beginverse{}
\ch{D}{}{}{}Nad stádem \ch{A}{}{}{}koní, \ch{Emi}{}{}{}podkovy \ch{G}{}{}{}zvoní, zvoní
\ch{D}{}{}{}Černý vůz \ch{A}{}{}{}vlečou \ch{Emi}{}{}{}a slzy \ch{G}{}{}{}tečou a já volám
\endverse{}

\beginverse{}
\ch{D}{}{}{}Tak neplač \ch{A}{}{}{}můj kamará\ch{Emi}{}{}{}de, náhoda je \ch{G}{}{}{}blbec, když krade
\ch{D}{}{}{}Je tuhý jak \ch{A}{}{}{}veka a řeka ho \ch{Emi}{}{}{}splaví, máme ho \ch{G}{}{}{}rádi
\endverse{}

\beginchorus{}
No tak \ch{C}{}{}{}có, tak \ch{G}{}{}{}có, tak \ch{A}{}{}{}có
\endchorus{}

\beginverse{}
\ch{D}{}{}{}Vždycky si \ch{A}{}{}{}přál, až bude \ch{Emi}{}{}{}popel, i s kyta\ch{G}{}{}{}rou
\ch{D}{}{}{}Vodou ať \ch{A}{}{}{}plavou, jen žádný \ch{Emi}{}{}{}hotel, s křížkem nad \ch{G}{}{}{}hlavou
\endverse{}

\beginverse{}
\ch{D}{}{}{}Až najdeš \ch{A}{}{}{}místo, kde je ten \ch{Emi}{}{}{}pramen a kámen co \ch{G}{}{}{}praská
\ch{D}{}{}{}Budeš mít \ch{A}{}{}{}jisto, patří sem \ch{Emi}{}{}{}popel a každá \ch{G}{}{}{}láska
\endverse{}

\repeatingchorus{}
\beginchorus{}
No tak \ch{C}{}{}{}có, tak \ch{G}{}{}{}có, tak \ch{A}{}{}{}có
\endchorus{}

\beginverse{}
\ch{D}{}{}{}Nad stádem \ch{A}{}{}{}koní, \ch{Emi}{}{}{}podkovy \ch{G}{}{}{}zvoní, zvoní
\ch{D}{}{}{}Černý vůz \ch{A}{}{}{}vlečou \ch{Emi}{}{}{}a slzy \ch{G}{}{}{}tečou a já šeptám
\ch{D}{}{}{}Vysyp ten \ch{A}{}{}{}popel \ch{Emi}{}{}{}kamará\ch{G}{}{}{}de, \ch{D}{}{}{}do bílé \ch{A}{}{}{}vody, \ch{Emi}{}{}{}vody \ch{G}{}{}{}
\ch{D}{}{}{}Vyhasnul \ch{A}{}{}{}kotel a náhoda \ch{Emi}{}{}{}je, štěstí od podko\ch{G}{}{}{}vy.
\endverse{}

\beginverse{}
\lrep{} \ch{D}{}{}{}Vysyp ten \ch{A}{}{}{}popel \ch{G}{}{}{}kamará\ch{D}{}{}{}de, do bílé \ch{A}{}{}{}vody, \ch{G}{}{}{}vo\ch{D}{}{}{}dy,
vyhasnul \ch{D}{}{}{}kotel a \ch{A}{}{}{}náhoda \ch{Emi}{}{}{}je, štěstí od podko\ch{G}{}{}{}vy. \rrep{}
\endverse{}

\endsong{}
\sclearpage{}
% chktex-file 8
% chktex-file 12
% chktex-file 18
\beginsong{Nagasaki Hirošima}[by={Mňága \& Žďorp}]
% \capo{0}
\transpose{0}

\beginchorus{}
\ch{C}{}{}{}Tramvají \ch{G}{}{}{}dvojkou \ch{F}{}{}{}jezdíval jsem \ch{G}{}{}{}do Žide\ch{C}{}{}{}nic, \ch{G}{}{}{} \ch{F}{}{}{} \ch{G}{}{}{}
\ch{C}{}{}{}z takový \ch{G}{}{}{}lásky \ch{F}{}{}{}většinou \ch{G}{}{}{}nezbyde \ch{Ami}{}{}{}nic,
\ch{F}{}{}{}z takový \ch{C}{}{}{}lásky \ch{F}{}{}{}jsou kruhy \ch{C}{}{}{}pod oči\ch{G}{}{}{}ma
a \ch{C}{}{}{}dvě spálený \ch{G}{}{}{}srdce -- \ch{F}{}{}{}Nagasaki, \ch{G}{}{}{}Hiroši\ch{C}{}{}{}ma. \ch{G}{}{}{} \ch{F}{}{}{} \ch{G}{}{}{}
\endchorus{}

\beginverse{}
\ch{C}{}{}{}Jsou jistý \ch{G}{}{}{}věci, \ch{F}{}{}{}co bych tesal \ch{G}{}{}{}do kame\ch{C}{}{}{}ne, \ch{G}{}{}{} \ch{F}{}{}{} \ch{G}{}{}{}
\ch{C}{}{}{}tam, kde je \ch{G}{}{}{}láska, \ch{F}{}{}{}tam je všechno \ch{G}{}{}{}dovole\ch{Ami}{}{}{}né,
\ch{F}{}{}{}a tam, kde \ch{C}{}{}{}není, \ch{F}{}{}{}tam mě to \ch{C}{}{}{}nezají\ch{G}{}{}{}má,
jó, \ch{C}{}{}{}dvě spálený \ch{G}{}{}{}srdce -- \ch{F}{}{}{}Nagasaki, \ch{G}{}{}{}Hiroši\ch{C}{}{}{}ma. \ch{G}{}{}{} \ch{F}{}{}{} \ch{G}{}{}{}
\endverse{}

\beginverse{}
\ch{C}{}{}{}Já nejsem \ch{G}{}{}{}svatej, \ch{F}{}{}{}ani ty \ch{G}{}{}{}nejsi sva\ch{C}{}{}{}tá, \ch{G}{}{}{} \ch{F}{}{}{} \ch{G}{}{}{}
\ch{C}{}{}{}jablka z \ch{G}{}{}{}ráje \ch{F}{}{}{}bejvala \ch{G}{}{}{}jedova\ch{Ami}{}{}{}tá,
jenže! -- \ch{F}{}{}{}hezky jsi \ch{C}{}{}{}hřála, \ch{F}{}{}{}když mi někdy \ch{C}{}{}{}byla zi\ch{G}{}{}{}ma,
jó, dvě \ch{C}{}{}{}spálený \ch{G}{}{}{}srdce -- \ch{F}{}{}{}Nagasaki, \ch{G}{}{}{}Hiroši\ch{C}{}{}{}ma. \ch{G}{}{}{} \ch{F}{}{}{} \ch{G}{}{}{}
\endverse{}

\beginchorus{}
\ch{C}{}{}{}Tramvají \ch{G}{}{}{}dvojkou \ch{F}{}{}{}jezdíval jsem \ch{G}{}{}{}do Žide\ch{C}{}{}{}nic, \ch{G}{}{}{} \ch{F}{}{}{} \ch{G}{}{}{}
\ch{C}{}{}{}z takový \ch{G}{}{}{}lásky \ch{F}{}{}{}většinou \ch{G}{}{}{}nezbyde \ch{Ami}{}{}{}nic,
\ch{F}{}{}{}z takový \ch{C}{}{}{}lásky \ch{F}{}{}{}jsou kruhy \ch{C}{}{}{}pod oči\ch{G}{}{}{}ma
a \ch{C}{}{}{}dvě spálený \ch{G}{}{}{}srdce -- \ch{F}{}{}{}Nagasaki, \ch{G}{}{}{}Hiroši\ch{C}{}{}{}ma. \ch{G}{}{}{} \ch{F}{}{}{} \ch{G}{}{}{}
a \ch{C}{}{}{}dvě spálený \ch{G}{}{}{}srdce -- \ch{F}{}{}{}Nagasaki, \ch{G}{}{}{}Hiroši\ch{C}{}{}{}ma. \ch{G}{}{}{} \ch{F}{}{}{} \ch{G}{}{}{}
a \ch{C}{}{}{}dvě spálený \ch{G}{}{}{}srdce -- \ch{F}{}{}{}Nagasaki, \ch{G}{}{}{}Hiroši\ch{C}{}{}{}ma. \ch{G}{}{}{} \ch{F}{}{}{} \ch{G}{}{}{}
a \ch{C}{}{}{}dvě spálený \ch{G}{}{}{}srdce -- \ch{F}{}{}{}Nagasaki, \ch{G}{}{}{}Hiroši\ch{C}{}{}{}ma. \ch{G}{}{}{} \ch{F}{}{}{} \ch{G}{}{}{}
\endchorus{}

\endsong{}
\sclearpage{}
\input{songs/oddilovy_zpevnik_II/src/Naj_naj+musician.tex}\sclearpage{}
\input{songs/oddilovy_zpevnik_II/src/Nápady_-_Dej_mi_víc_své_lásky-Olympic+musician.tex}\sclearpage{}
\input{songs/oddilovy_zpevnik_II/src/Niagára+musician.tex}\sclearpage{}
% chktex-file 8
% chktex-file 12
% chktex-file 18
\beginsong{Omnia vincit Amor}[by={Klíč}]
% \capo{0}
\transpose{0}

\beginverse*{}
\ch{F}{}{}{} \ch{C}{}{}{} \ch{Dmi}{}{}{} \ch{Ami}{}{}{} \ch{B{\flt}}{}{}{} \ch{C}{}{}{} \ch{Dmi}{}{}{} \ch{F}{}{}{} \ch{C}{}{}{} \ch{Dmi}{}{}{} \ch{Ami}{}{}{} \ch{B{\flt}}{}{}{} \ch{C}{}{}{} \ch{Dmi}{}{}{}
\endverse{}

\beginverse{}
\ch{Dmi}{}{}{}Šel pocestný kol \ch{C}{}{}{}hospodských \ch{Dmi}{}{}{}zdí,
\ch{F}{}{}{}přisedl k nám a \ch{C}{}{}{}lokálem \ch{F}{}{}{}zní
\ch{Gmi}{}{}{}pozdrav jak svaté \ch{F}{}{}{}přikázá\ch{C}{}{}{}ní:
\ch{Dmi}{}{}{}omnia \ch{C}{}{}{}vincit \ch{Dmi}{}{}{}Amor.
\endverse{}

\beginverse{}
"\ch{Dmi}{}{}{}Hej, šenkýři, \ch{C}{}{}{}dej plný \ch{Dmi}{}{}{}džbán,
\ch{F}{}{}{}ať chasa ví, kdo k \ch{C}{}{}{}stolu je \ch{F}{}{}{}zván,
\ch{Gmi}{}{}{}se mnou ať zpívá, \ch{F}{}{}{}kdo za své \ch{C}{}{}{}vzal
\ch{Dmi}{}{}{}omnia \ch{C}{}{}{}vincit \ch{Dmi}{}{}{}Amor."
\endverse{}

\beginchorus{}
Zlaťák \ch{F}{}{}{}pálí, \ch{C}{}{}{}nesleví \ch{Dmi}{}{}{}nic,
štěstí v \ch{F}{}{}{}lásce \ch{C}{}{}{}znamená \ch{F}{}{}{}víc,
všechny \ch{Gmi}{}{}{}pány \ch{F}{}{}{}ať vezme \ch{C}{}{}{}ďas! \ch{A7}{}{}{}
\ch{Dmi}{}{}{}Omnia \ch{C}{}{}{}vincit \ch{Dmi}{}{}{}Amor.
\endchorus{}

\beginverse{}
"\ch{Dmi}{}{}{}Já viděl zemi \ch{C}{}{}{}válkou se \ch{Dmi}{}{}{}chvět,
\ch{F}{}{}{}musel se bít a \ch{C}{}{}{}nenávi\ch{F}{}{}{}dět,
v \ch{Gmi}{}{}{}plamenech pálit \ch{F}{}{}{}prosby a \ch{C}{}{}{}pláč,
\ch{Dmi}{}{}{}omnia \ch{C}{}{}{}vincit \ch{Dmi}{}{}{}Amor.
\endverse{}

\beginverse{}
\ch{Dmi}{}{}{}Zlý trubky troubí, \ch{C}{}{}{}vítězí \ch{Dmi}{}{}{}zášť,
\ch{F}{}{}{}nad lidskou láskou \ch{C}{}{}{}roztáhli \ch{F}{}{}{}plášť,
\ch{Gmi}{}{}{}vtom kdosi krví \ch{F}{}{}{}napsal ten \ch{C}{}{}{}vzkaz:
\ch{Dmi}{}{}{}omnia \ch{C}{}{}{}vincit \ch{Dmi}{}{}{}Amor."
\endverse{}

\beginchorus{}
Zlaťák \ch{F}{}{}{}pálí, \ch{C}{}{}{}nesleví \ch{Dmi}{}{}{}nic,
štěstí v \ch{F}{}{}{}lásce \ch{C}{}{}{}znamená \ch{F}{}{}{}víc,
všechny \ch{Gmi}{}{}{}pány \ch{F}{}{}{}ať vezme \ch{C}{}{}{}ďas! \ch{A7}{}{}{}
\ch{Dmi}{}{}{}Omnia \ch{C}{}{}{}vincit \ch{Dmi}{}{}{}Amor.
\endchorus{}

\beginverse{}
"\ch{Dmi}{}{}{}Já prošel každou z \ch{C}{}{}{}nejdelších \ch{Dmi}{}{}{}cest,
\ch{F}{}{}{}všude se ptal, co \ch{C}{}{}{}značí ta \ch{F}{}{}{}zvěst,
\ch{Gmi}{}{}{}až řekl moudrý:"\ch{F}{}{}{}Pochopíš \ch{C}{}{}{}sám
\ch{Dmi}{}{}{}omnia \ch{C}{}{}{}vincit \ch{Dmi}{}{}{}Amor -- všechno přemáhá láska."
\endverse{}

\beginchorus{}
Zlaťák \ch{F}{}{}{}pálí, \ch{C}{}{}{}nesleví \ch{Dmi}{}{}{}nic,
štěstí v \ch{F}{}{}{}lásce \ch{C}{}{}{}znamená \ch{F}{}{}{}víc,
všechny \ch{Gmi}{}{}{}pány \ch{F}{}{}{}ať vezme \ch{C}{}{}{}ďas! \ch{A7}{}{}{}
\ch{Dmi}{}{}{}Omnia \ch{C}{}{}{}vincit \ch{Dmi}{}{}{}Amor.
\endchorus{}

\beginverse{}
"\ch{Dmi}{}{}{}Teď s novou vírou \ch{C}{}{}{}obcházím \ch{Dmi}{}{}{}svět,
\ch{F}{}{}{}má hlava zšedla \ch{C}{}{}{}pod tíhou \ch{F}{}{}{}let,
\ch{Gmi}{}{}{}každého zdravím \ch{F}{}{}{}větou všech \ch{C}{}{}{}vět:
\ch{Dmi}{}{}{}omnia \ch{C}{}{}{}vincit \ch{Dmi}{}{}{}Amor, omnia \ch{C}{}{}{}vincit \ch{Dmi}{}{}{}Amor \ldots{}"
\endverse{}

\endsong{}
\sclearpage{}
% chktex-file 8
% chktex-file 12
% chktex-file 18
\beginsong{Pocity}[by={Tomáš Klus}]
% \capo{0}
\transpose{0}

\beginverse{}
\ch{C}{}{}{}Z posledních poci\ch{G}{}{}{}tů \ch{Ami}{}{}{}poskládám ještě jednou \ch{F}{}{}{}úžasnou chvíli
\ch{C}{}{}{}Je to tím, že jsi \ch{G}{}{}{}tu, \ch{Ami}{}{}{}možná tím, že \ch{F}{}{}{}kdysi jsme byli Ty a
\ch{C}{}{}{}Já my dva, dvě \ch{G}{}{}{}nahý těla, tak \ch{Ami}{}{}{}neříkej, že \ch{F}{}{}{}jinak si to chtěla,
tak \ch{C}{}{}{}neříkej, \ch{G}{}{}{}neříkej, \ch{Ami}{}{}{}neříkej mi \ch{F}{}{}{}nic.
\endverse{}

\beginverse{}
\ch{C}{}{}{}Stala ses do no\ch{G}{}{}{}ci \ch{Ami}{}{}{}z ničeho nic moje \ch{F}{}{}{}platonická láska
\ch{C}{}{}{}unaven, bezmoc\ch{G}{}{}{}ný, \ch{Ami}{}{}{}usínám vedle Tebe, \ch{F}{}{}{}něco ve mně praská.
\ch{C}{}{}{}A ranní probuze\ch{G}{}{}{}ní a slova o štěs\ch{Ami}{}{}{}tí, neboj se to nic ne\ch{F}{}{}{}ní
\ch{C}{}{}{}Pohled a okouzle\ch{G}{}{}{}ní, a prázdný náměstí \ch{Ami}{}{}{}na znamení. \ch{F}{}{}{}
\endverse{}

\beginchorus{}
Jenže ty \ch{C}{}{}{}neslyšíš,\ch{G}{}{}{} jenže ty \ch{Ami}{}{}{}neposloucháš,\ch{F}{}{}{}
snad ani \ch{C}{}{}{}nevidíš,\ch{G}{}{}{} nebo spíš \ch{Ami}{}{}{}nechceš vidět\ch{F}{}{}{}
a druhejm \ch{C}{}{}{}závidíš\ch{G}{}{}{} a v očích \ch{Ami}{}{}{}kapky slaný vody\ch{F}{}{}{}
zkus změnu, \ch{C}{}{}{}uvidíš.\ch{G}{}{}{} Pak \ch{Ami}{}{}{}vítej do svobody\ch{F}{}{}{}
\endchorus{}

\beginverse{}
\ch{C}{}{}{}Jsi anděl, netu\ch{G}{}{}{}šíš. \ch{Ami}{}{}{}Anděl, co ze strachu \ch{F}{}{}{}mu utrhali křídla.
\ch{C}{}{}{}A až to ucí\ch{G}{}{}{}tíš, \ch{Ami}{}{}{}zkus kašlat na pravid\ch{F}{}{}{}la.
\ch{C}{}{}{}Říkej si o mně, co \ch{G}{}{}{}chceš. \ch{Ami}{}{}{}Já jsem byl \ch{F}{}{}{}odjakživa blázen.
\ch{C}{}{}{}Nevím, co nechá\ch{G}{}{}{}peš, ale \ch{Ami}{}{}{}vrať se na zem.\ch{F}{}{}{}
\endverse{}

\beginchorus{}
Jenže ty \ch{C}{}{}{}neslyšíš,\ch{G}{}{}{} jenže ty \ch{Ami}{}{}{}neposloucháš,\ch{F}{}{}{}
snad ani \ch{C}{}{}{}nevidíš,\ch{G}{}{}{} nebo spíš \ch{Ami}{}{}{}nechceš vidět\ch{F}{}{}{}
a druhejm \ch{C}{}{}{}závidíš\ch{G}{}{}{} a v očích \ch{Ami}{}{}{}kapky slaný vody\ch{F}{}{}{}
zkus změnu, \ch{C}{}{}{}uvidíš.\ch{G}{}{}{} Pak \ch{Ami}{}{}{}vítej do svobody\ch{F}{}{}{}
\endchorus{}

\beginverse{}
\ch{C}{}{}{}Z posledních poci\ch{G}{}{}{}tů \ch{Ami}{}{}{}poskládám ještě jednou \ch{F}{}{}{}úžasnou chvíli
\ch{C}{}{}{}Je to tím, že jsi \ch{G}{}{}{}tu, \ch{Ami}{}{}{}možná tím, že \ch{F}{}{}{}kdysi jsme byli Ty a
\ch{C}{}{}{}Já my dva, dvě \ch{G}{}{}{}nahý těla, tak \ch{Ami}{}{}{}neříkej, že \ch{F}{}{}{}jinak si to chtěla,
tak \ch{C}{}{}{}neříkej, \ch{G}{}{}{}neříkej, \ch{Ami}{}{}{}neříkej mi \ch{F}{}{}{}nic.
\endverse{}

\endsong{}
\sclearpage{}
% chktex-file 8
% chktex-file 12
% chktex-file 18
\beginsong{Podvod}[by={Nedvědi}]
% \capo{0}
\transpose{0}

\beginverse{}
\ch{Emi}{}{}{}Na dlani jednu z tvých řas, do tmy se \ch{Ami}{}{}{}koukám,
\ch{D}{}{}{}hraju si písničky tvý, co jsem ti \ch{G}{}{}{}psal,
je skoro \ch{Ami}{}{}{}půlnoc a z kostela \ch{C}{}{}{}zvon mi noc připo\ch{Emi}{}{}{}míná,
půjdu se \ch{Ami}{}{}{}mejt a pozhasínám, co bude \ch{B7}{}{}{}dál?
\endverse{}

\beginverse{}
\ch{Emi}{}{}{}Pod polštář dopisů pár, co poslalas, \ch{Ami}{}{}{}dávám,
\ch{D}{}{}{}píšeš, že ráda mě máš a trápí tě \ch{G}{}{}{}stesk,
je skoro \ch{Ami}{}{}{}půlnoc a z kostela \ch{C}{}{}{}zvon mi noc připo\ch{Emi}{}{}{}míná,
půjdu se \ch{Ami}{}{}{}mejt a pozhasínám, co bude \ch{B7}{}{}{}dál?
\endverse{}

\beginchorus{}
Chtěl jsem to \ch{Ami}{}{}{}ráno, kdy naposled snídal jsem s tebou,
ti \ch{Emi}{}{}{}říct, že už ti nezavolám,
pro jednu pitomou \ch{Ami}{}{}{}holku, pro pár nocí \ch{D}{}{}{}touhy
podved jsem \ch{G}{}{}{}všechno, o čem doma si \ch{B7}{}{}{}sníš,
teď je mi to \ch{Emi}{}{}{}líto.
\endchorus{}

\beginverse{}
\ch{Emi}{}{}{}Kolikrát člověk může mít rád tak opravdu z \ch{Ami}{}{}{}lásky,
\ch{D}{}{}{}dvakrát či třikrát -- to ne, i jednou je \ch{G}{}{}{}dost,
je skoro \ch{Ami}{}{}{}půlnoc a z kostela \ch{C}{}{}{}zvon mi noc připo\ch{Emi}{}{}{}míná,
půjdu se \ch{Ami}{}{}{}mejt a pozhasínám, co bude \ch{B7}{}{}{}dál?
\endverse{}

\beginchorus{}
Chtěl jsem to \ch{Ami}{}{}{}ráno, kdy naposled snídal jsem s tebou,
ti \ch{Emi}{}{}{}říct, že už ti nezavolám,
pro jednu pitomou \ch{Ami}{}{}{}holku, pro pár nocí \ch{D}{}{}{}touhy
podved jsem \ch{G}{}{}{}všechno, o čem doma si \ch{B7}{}{}{}sníš,
teď je mi to \ch{Emi}{}{}{}líto.
\endchorus{}

\beginchorus{}
Chtěl jsem to \ch{Ami}{}{}{}ráno, kdy naposled snídal jsem s tebou,
ti \ch{Emi}{}{}{}říct, že už ti nezavolám,
pro jednu pitomou \ch{Ami}{}{}{}holku, pro pár nocí \ch{D}{}{}{}touhy
podved jsem \ch{G}{}{}{}všechno, o čem doma si \ch{B7}{}{}{}sníš,
teď je mi to \ch{Emi}{}{}{}líto.
\endchorus{}

\endsong{}
\sclearpage{}
\input{songs/oddilovy_zpevnik_II/src/Pohádky_je_konec-Lucie_Bílá+musician.tex}\sclearpage{}
\input{songs/oddilovy_zpevnik_II/src/Pojď_sem+musician.tex}\sclearpage{}
\input{songs/oddilovy_zpevnik_II/src/Poštorenská_kapela+musician.tex}\sclearpage{}
% chktex-file 8
% chktex-file 12
% chktex-file 18
\beginsong{Ráda se miluje}[by={Karel Plíhal}]
% \capo{0}
\transpose{0}

\beginchorus{}
\ch{Bmi}{}{}{}Ráda se miluje, \ch{A}{}{}{}ráda \ch{D}{}{}{}jí, \ch{G}{}{}{}ráda si \ch{F{\shrp}mi}{}{}{}jenom tak \ch{Bmi}{}{}{}zpívá,
\ch{Bmi}{}{}{}vrabci se na plotě \ch{A}{}{}{}háda\ch{D}{}{}{}jí, \ch{G}{}{}{}kolik že \ch{F{\shrp}mi}{}{}{}času jí \ch{Bmi}{}{}{}zbývá.
\endchorus{}

\beginverse{}
\ch{G}{}{}{}Než vítr dostrká k \ch{D}{}{}{}útesu \ch{G}{}{}{}tu její legrační \ch{D}{}{}{}bár\ch{F{\shrp}}{}{}{}ku
a \ch{Bmi}{}{}{}Pámbu si ve svým \ch{A}{}{}{}note\ch{D}{}{}{}su \ch{G}{}{}{}udělá \ch{F{\shrp}mi}{}{}{}jen další \ch{Bmi}{}{}{}čárku.
\endverse{}

\repeatingchorus{}
\beginchorus{}
\ch{Bmi}{}{}{}Ráda se miluje, \ch{A}{}{}{}ráda \ch{D}{}{}{}jí, \ch{G}{}{}{}ráda si \ch{F{\shrp}mi}{}{}{}jenom tak \ch{Bmi}{}{}{}zpívá,
\ch{Bmi}{}{}{}vrabci se na plotě \ch{A}{}{}{}háda\ch{D}{}{}{}jí, \ch{G}{}{}{}kolik že \ch{F{\shrp}mi}{}{}{}času jí \ch{Bmi}{}{}{}zbývá.
\endchorus{}

\beginverse{}
\ch{G}{}{}{}Psáno je v nebeské \ch{D}{}{}{}režii, \ch{G}{}{}{}a to hned na první \ch{D}{}{}{}strán\ch{F{\shrp}}{}{}{}ce,
\ch{Bmi}{}{}{}že naše duše nás \ch{A}{}{}{}přeži\ch{D}{}{}{}jí v \ch{G}{}{}{}jinačí \ch{F{\shrp}mi}{}{}{}tělesný \ch{Bmi}{}{}{}schránce.
\endverse{}

\repeatingchorus{}
\beginchorus{}
\ch{Bmi}{}{}{}Ráda se miluje, \ch{A}{}{}{}ráda \ch{D}{}{}{}jí, \ch{G}{}{}{}ráda si \ch{F{\shrp}mi}{}{}{}jenom tak \ch{Bmi}{}{}{}zpívá,
\ch{Bmi}{}{}{}vrabci se na plotě \ch{A}{}{}{}háda\ch{D}{}{}{}jí, \ch{G}{}{}{}kolik že \ch{F{\shrp}mi}{}{}{}času jí \ch{Bmi}{}{}{}zbývá.
\endchorus{}

\beginverse{}
\ch{G}{}{}{}Úplně na konci \ch{D}{}{}{}paseky, \ch{G}{}{}{}tam, kde se ozvěna \ch{D}{}{}{}tříš\ch{F{\shrp}}{}{}{}tí,
\ch{Bmi}{}{}{}sedí šnek ve snacku \ch{A}{}{}{}pro šne\ch{D}{}{}{}ky -- \ch{G}{}{}{}snad její \ch{F{\shrp}mi}{}{}{}podoba \ch{Bmi}{}{}{}příští.
\endverse{}

\repeatingchorus{}
\beginchorus{}
\ch{Bmi}{}{}{}Ráda se miluje, \ch{A}{}{}{}ráda \ch{D}{}{}{}jí, \ch{G}{}{}{}ráda si \ch{F{\shrp}mi}{}{}{}jenom tak \ch{Bmi}{}{}{}zpívá,
\ch{Bmi}{}{}{}vrabci se na plotě \ch{A}{}{}{}háda\ch{D}{}{}{}jí, \ch{G}{}{}{}kolik že \ch{F{\shrp}mi}{}{}{}času jí \ch{Bmi}{}{}{}zbývá.
\endchorus{}

\endsong{}
\sclearpage{}
\input{songs/oddilovy_zpevnik_II/src/Rosa_na_kolejích-Wabi_Daněk+musician.tex}\sclearpage{}
\input{songs/oddilovy_zpevnik_II/src/Rozárka-Divokej_Bill+musician.tex}\sclearpage{}
\input{songs/oddilovy_zpevnik_II/src/Růže_z_papíru-Nedvědi+musician.tex}\sclearpage{}
\input{songs/oddilovy_zpevnik_II/src/Sára-Traband+musician.tex}\sclearpage{}
\input{songs/oddilovy_zpevnik_II/src/Síla_starejch_vín-Škwor+musician.tex}\sclearpage{}
% chktex-file 8
% chktex-file 12
% chktex-file 18
\beginsong{Snídaně v trávě}[by={Michal Tučný}]
% \capo{0}
\transpose{0}

\beginverse{}
\ch{G}{}{}{}Já nejsem \ch{Emi}{}{}{}zdejší, \ch{C}{}{}{}jen projíž\ch{D}{}{}{}dím,
\ch{G}{}{}{}jsou movi\ch{Emi}{}{}{}tější, \ch{C}{}{}{}já ti nabí\ch{D}{}{}{}zím:
\endverse{}

\beginchorus{}
Snídani v \ch{G}{}{}{}trávě, teď právě \ch{D}{}{}{}prostřenou,
loukou ráno \ch{G}{}{}{}vzbuzenou, napůl v snách
\ch{C}{}{}{}hledám po kapsách, čím \ch{G}{}{}{}zaplatím
dřív, než se \ch{D}{}{}{}vytratím, snídani v \ch{G}{}{}{}trávě.
\endchorus{}

\beginverse{}
\ch{G}{}{}{}Dál podle \ch{Emi}{}{}{}pražců, \ch{C}{}{}{}tam je můj \ch{D}{}{}{}cíl,
\ch{G}{}{}{}dám na každej \ch{Emi}{}{}{}nápis, \ch{C}{}{}{}že jsem tu \ch{D}{}{}{}žil.
\endverse{}

\repeatingchorus{}
\beginchorus{}
Snídani v \ch{G}{}{}{}trávě, teď právě \ch{D}{}{}{}prostřenou,
loukou ráno \ch{G}{}{}{}vzbuzenou, napůl v snách
\ch{C}{}{}{}hledám po kapsách, čím \ch{G}{}{}{}zaplatím
dřív, než se \ch{D}{}{}{}vytratím, snídani v \ch{G}{}{}{}trávě.
\endchorus{}

\beginverse*{}
\ch{C}{}{}{}Kdybych měl plátno \ch{G}{}{}{}a zlatý rám,
\ch{C}{}{}{}tak se do malování \ch{B}{}{}{}dám.
\endverse{}

\beginverse{}
\ch{G}{}{}{}Jsou povola\ch{Emi}{}{}{}nější, \ch{C}{}{}{}jak jsem si \ch{D}{}{}{}všim',
\endverse{}

\beginverse{}
\ch{G}{}{}{}nejsem z\ch{Emi}{}{}{}dejší, \ch{C}{}{}{}tak jenom \ch{D}{}{}{}sním.
\endverse{}

\repeatingchorus{}
\beginchorus{}
Snídani v \ch{G}{}{}{}trávě, teď právě \ch{D}{}{}{}prostřenou,
loukou ráno \ch{G}{}{}{}vzbuzenou, napůl v snách
\ch{C}{}{}{}hledám po kapsách, čím \ch{G}{}{}{}zaplatím
dřív, než se \ch{D}{}{}{}vytratím, snídani v \ch{G}{}{}{}trávě.
\endchorus{}

\endsong{}
\sclearpage{}
% chktex-file 8
% chktex-file 12
% chktex-file 18
\beginsong{Strč prst skrz krk}[by={Bobři}]
% \capo{0}
\transpose{0}

\beginchorus{}
\ch{G}{}{}{}Strč prst skrz krk zřetelně nám přeříkej,
\ch{D}{}{}{}strč prst skrz krk zřetelně nám přeříkej,
\ch{G}{}{}{}strč prst skrz krk zřetelně nám přeříkej,
\ch{D}{}{}{}bude z tebe zpěvák velikej, \ch{G}{}{}{}bájo.
\endchorus{}

\beginverse{}
Na cvičišti čtyři svišti piští,
těším se na lekci příští, bájo.
\endverse{}

\beginverse{}
Dalajláma v lomu láme slámu,
já si na tom jazyk zlámu, bájo.
\endverse{}

\beginverse{}
Jak Julie olejuje koleje,
smrtelný pot ze mě leje, bájo.
\endverse{}

\beginverse{}
Teď už umím "kaplan plakal v kapli",
hlasivky se mi už naply, bájo.
\endverse{}

\beginverse{}
Od poklopu ku poklopu Kyklop kouli koulí,
už mám z toho na jazyku bouli, bájo.
\endverse{}

\beginverse{}
Skákalo psisko přes příkopisko
hop sem, hop tam, v rytmu disco, bájo.
\endverse{}

\beginverse{}
Na klavír hrála královna Klára,
a přitom se na mě smála, bájo.
\endverse{}

\beginverse{}
Tato teta to je této tety teta,
je to bída tahle věta, bájo.
\endverse{}

\beginverse{}
10: Seno snese se se sesečené louky,
není to nic pro samouky, bájo.
\endverse{}

\endsong{}
\sclearpage{}
\input{songs/oddilovy_zpevnik_II/src/Stýskání-Bratři_Ebenové+musician.tex}\sclearpage{}
\input{songs/oddilovy_zpevnik_II/src/Tak_už_mi_má_holka_mává-Wabi_Ryvola+musician.tex}\sclearpage{}
% chktex-file 8
% chktex-file 12
% chktex-file 18
\beginsong{To ta Helpa}
% \capo{0}
\transpose{0}

\beginverse{}
\ch{Dmi}{}{}{}To ta \ch{G}{}{}{}Helpa, \ch{Dmi}{}{}{}to ta \ch{G}{}{}{}Helpa, \ch{Dmi}{}{}{}to je \ch{A7}{}{}{}pekné \ch{Dmi}{}{}{}mesto
\ch{Dmi}{}{}{}a v tej \ch{G}{}{}{}Helpe, \ch{Dmi}{}{}{}a v tej \ch{G}{}{}{}Helpe \ch{Dmi}{}{}{}švarných \ch{A7}{}{}{}chlapcov \ch{Dmi}{}{}{}je sto.
\ch{B{\flt}}{}{}{}Koho \ch{C7}{}{}{}je sto, \ch{F}{}{}{}toho je sto, \ch{C}{}{}{}nie po mojej \ch{F}{}{}{}vó\ch{A7}{}{}{}li,
\ch{Dmi}{}{}{}len za \ch{G}{}{}{}jednym, \ch{Dmi}{}{}{}len za \ch{G}{}{}{}jednym \ch{Dmi}{}{}{}srdieč\ch{A7}{}{}{}ko ma \ch{Dmi}{}{}{}boli
\ch{B{\flt}}{}{}{}Koho \ch{C7}{}{}{}je sto, \ch{F}{}{}{}toho je sto, \ch{C}{}{}{}nie po mojej \ch{F}{}{}{}vó\ch{A7}{}{}{}li,
\ch{Dmi}{}{}{}len za \ch{G}{}{}{}jednym, \ch{Dmi}{}{}{}len za \ch{G}{}{}{}jednym \ch{Dmi}{}{}{}srdieč\ch{A7}{}{}{}ko ma \ch{Dmi}{}{}{}boli
\endverse{}

\beginverse{}
\ch{Dmi}{}{}{}Za Ja\ch{G}{}{}{}níčkom, \ch{Dmi}{}{}{}za Pa\ch{G}{}{}{}líčkom \ch{Dmi}{}{}{}krok by \ch{A7}{}{}{}něspra\ch{Dmi}{}{}{}vila,
\ch{Dmi}{}{}{}za Ďu\ch{G}{}{}{}ríčkom, \ch{Dmi}{}{}{}za Mi\ch{G}{}{}{}šíčkom \ch{Dmi}{}{}{}Dunaj \ch{A7}{}{}{}presko\ch{Dmi}{}{}{}čila.
\ch{B{\flt}}{}{}{}Dunaj, \ch{C7}{}{}{}Dunaj, \ch{F}{}{}{}Dunaj, Dunaj, \ch{C}{}{}{}aj to širé \ch{F}{}{}{}po\ch{A7}{}{}{}le,
\ch{Dmi}{}{}{}len za \ch{G}{}{}{}jedním, \ch{Dmi}{}{}{}len za \ch{G}{}{}{}jedním, \ch{Dmi}{}{}{}poče\ch{A7}{}{}{}šenie \ch{Dmi}{}{}{}moje.
\ch{B{\flt}}{}{}{}Dunaj, \ch{C7}{}{}{}Dunaj, \ch{F}{}{}{}Dunaj, Dunaj, \ch{C}{}{}{}aj to širé \ch{F}{}{}{}po\ch{A7}{}{}{}le,
\ch{Dmi}{}{}{}len za \ch{G}{}{}{}jedním, \ch{Dmi}{}{}{}len za \ch{G}{}{}{}jedním, \ch{Dmi}{}{}{}poče\ch{A7}{}{}{}šenie \ch{Dmi}{}{}{}moje.
\endverse{}

\endsong{}
\sclearpage{}
\input{songs/oddilovy_zpevnik_II/src/Tu_kytaru_jsem_koupil_kvůli_tobě-Václav_Neckář+musician.tex}\sclearpage{}
\input{songs/oddilovy_zpevnik_II/src/Tulácké_blues-Pacifik+musician.tex}\sclearpage{}
\input{songs/oddilovy_zpevnik_II/src/U_nás_na_severu-Jaromír_Nohavica+musician.tex}\sclearpage{}
\input{songs/oddilovy_zpevnik_II/src/Už_to_nenapravím-Samson+musician.tex}\sclearpage{}
\input{songs/oddilovy_zpevnik_II/src/Víno-Chinaski+musician.tex}\sclearpage{}
\input{songs/oddilovy_zpevnik_II/src/Vodácká_holka-Hop_Trop+musician.tex}\sclearpage{}
% chktex-file 8
% chktex-file 12
% chktex-file 18
\beginsong{Všichni jsou už v Mexiku}[by={Michal Tučný}]
% \capo{0}
\transpose{0}

\beginverse{}
\ch{D}{}{}{}Kde jen jsou a kde jsem já, \ch{A}{}{}{}proč všechno končí, co začíná,
\ch{Bmi}{}{}{}smůla se na mě lepí \ch{E}{}{}{}ve zvyku \ch{A}{}{}{}a všichni jsou už \ch{A7}{}{}{}v Mexiku.
\endverse{}

\beginchorus{}
\ch{D}{}{}{}Všichni jsou už v Mexiku, \ch{A}{}{}{}Buenos Dias já taky jdu,
\ch{Bmi}{}{}{}aspoň si poslechnu pěknou \ch{E}{}{}{}muziku, \ch{A}{}{}{}co se hraje \ch{A7}{}{}{}v Mexiku.
\endchorus{}

\beginverse{}
\ch{D}{}{}{}Kde je můj kůň, co měl mě nést \ch{A}{}{}{}blátem měst a prachem cest,
\ch{Bmi}{}{}{}dělal, že se ho to už \ch{E}{}{}{}netýká, \ch{A}{}{}{}asi šel do \ch{A7}{}{}{}Mexika.
\ch{D}{}{}{}Kde je ta holka, co měl jsem rád, \ch{A}{}{}{}už tu na rohu měla stát,
\ch{Bmi}{}{}{}s pihou na svým krásným \ch{E}{}{}{}nosíku \ch{A}{}{}{}už se asi nosí někde \ch{A7}{}{}{}v Mexiku
\endverse{}

\beginchorus{}
\ch{D}{}{}{}Všichni jsou už v Mexiku, \ch{A}{}{}{}Buenos Dias já taky jdu,
\ch{Bmi}{}{}{}aspoň si poslechnu pěknou \ch{E}{}{}{}muziku, \ch{A}{}{}{}co se hraje \ch{A7}{}{}{}v Mexiku.
\endchorus{}

\beginverse{}
\ch{D}{}{}{}Kde je můj pes, co mě hlídat měl, \ch{A}{}{}{}že nezaštěkal, kam jen šel
\ch{Bmi}{}{}{}bez obojku a bez \ch{E}{}{}{}košíku \ch{A}{}{}{}už se asi toulá v \ch{A7}{}{}{}Mexiku.
\ch{D}{}{}{}Kde je ten náš slavnej gang, \ch{A}{}{}{}postrach salonů a postrach bank
\ch{Bmi}{}{}{}nechali mi všechno pěkně \ch{E}{}{}{}na triku \ch{A}{}{}{}a už jsou někde dole v \ch{A7}{}{}{}Mexiku
\endverse{}

\beginchorus{}
\ch{D}{}{}{}Všichni jsou už v Mexiku, \ch{A}{}{}{}Buenos Dias já taky jdu,
\ch{Bmi}{}{}{}aspoň si poslechnu pěknou \ch{E}{}{}{}muziku, \ch{A}{}{}{}co se hraje \ch{A7}{}{}{}v Mexiku.
\endchorus{}

\beginverse{}
\ch{D}{}{}{}Kam šel barman, teď tu stál \ch{A}{}{}{}a chlapů houf, co s nima hrál,
\ch{Bmi}{}{}{}jen zatykač s mojí fotkou \ch{E}{}{}{}visí tu \ch{A}{}{}{}a ostatní jsou \ch{A7}{}{}{}v Mexiku.
\ch{D}{}{}{}Kde je ten chlápek, co tu spal \ch{A}{}{}{}a ta kočka, co jí nalejval,
\ch{Bmi}{}{}{}šerif právě bere v baru \ch{E}{}{}{}za kliku, \ch{A}{}{}{}asi jsem měl bejt už taky \ch{A7}{}{}{}v Mexiku.
\endverse{}

\beginchorus{}
\ch{D}{}{}{}Všichni jsou už v Mexiku, \ch{A}{}{}{}Buenos Dias já taky jdu,
\ch{Bmi}{}{}{}aspoň si poslechnu pěknou \ch{E}{}{}{}muziku, \ch{A}{}{}{}co se hraje \ch{A7}{}{}{}v Mexiku.
\endchorus{}

\endsong{}
\sclearpage{}
\input{songs/oddilovy_zpevnik_II/src/Zabili_zabili+musician.tex}\sclearpage{}
% chktex-file 8
% chktex-file 12
% chktex-file 18
\beginsong{Zahrada ticha}[by={Jakub Smolík}]
% \capo{0}
\transpose{0}

\beginverse{}
Je tam brána zdobe\ch{D}{}{}{}ná, cestu oteví\ch{Em}{}{}{}rá,
zahradu zele\ch{C}{}{}{}nou,\ch{G}{}{}{} všechno připomí\ch{D}{}{}{}ná.
\endverse{}

\beginverse{}
Jako dým závo\ch{D}{}{}{}jů, mlhou upřede\ch{Emi}{}{}{}ných,
vstupuješ do ti\ch{C}{}{}{}cha,\ch{G}{}{}{} cestou vyvole\ch{D}{}{}{}ných.
\endverse{}

\beginverse{}
Je to březový \ch{D}{}{}{}háj, je to borový \ch{Emi}{}{}{}les,
je to anglický \ch{C}{}{}{}park,\ch{G}{}{}{} je to hluboký \ch{D}{}{}{}vřes.
\endverse{}

\beginverse{}
Je to samota \ch{D}{}{}{}dnů, kdy jsi pomalu \ch{Emi}{}{}{}zrál,
v zahradě zele\ch{C}{}{}{}ný,\ch{G}{}{}{} kde sis za dětství \ch{D}{}{}{}hrál.
\endverse{}

\beginverse{}
Kolik chceš, tolik \ch{D}{}{}{}máš, očí otevře\ch{Emi}{}{}{}ných,
tam venku za bra\ch{C}{}{}{}nou,\ch{G}{}{}{} leží studený \ch{D}{}{}{}sníh.
\endverse{}

\beginverse{}
Z počátku usly\ch{D}{}{}{}šíš, vítr a ptačí \ch{Emi}{}{}{}hlas,
v zahradě zele\ch{C}{}{}{}ný,\ch{G}{}{}{} přejdou do ticha \ch{D}{}{}{}zas.
\endverse{}

\beginverse{}
Světlo připomí\ch{D}{}{}{}ná, rána slunečných \ch{Emi}{}{}{}dnů,
v zahradě zele\ch{C}{}{}{}ný,\ch{G}{}{}{} v zahradě beze \ch{D}{}{}{}snů.
\endverse{}

\beginverse{}
Uprostřed závra\ch{D}{}{}{}tí, sluncem prosvíce\ch{Emi}{}{}{}ných,
vstupuješ do ti\ch{C}{}{}{}cha,\ch{G}{}{}{} cestou vyvole\ch{D}{}{}{}ných.
\endverse{}

\endsong{}
\sclearpage{}
\input{songs/oddilovy_zpevnik_II/src/Za_svou_pravdou_stát-Spirituál_kvintet+musician.tex}\sclearpage{}

  }
\end{songs}

\end{document}
